\scanpage{58}
% archive sheet number: 57

3.1.84

Transport parallèle le long d'un demi-tour, \deleted{\unsure{au}} le long d'une position relative \deleted{singulière de position} supersingulière

\[ \boxed{\lambda(t)=t-(n-1)} \]

\uline{NB} Le pgcd de $n-1$ et de $2n$ est 1 si $n$ pair, 2 si $n$ impair.

% slanted
\margin{\unsure{Toute|Tout} la suite \unsure{couillon|couill.} !}

\deleted{Donc on voit que l'image de
\[ \Pi_1(\mathcal X)\longrightarrow\deleted{\mathbb Z}\mathbb D_{2n} \]
a une image qui contient $\mathbb D_{2n}^{+}\inserted{\simeq\mathbb Z/n\mathbb Z}$ si $n$ pair (donc cette image est d'indice 1 ou 2) \deleted{ou bien elle} et elle contient $\mathbb Z/n\mathbb Z$ si $n$ impair (donc elle est alors d'indice 1, 2 ou 4)

\deleted{Cela ne nous dit rien sur les compos} Si on prend le dédoublage $\widetilde{\mathcal X}$ de $\mathcal X$ suivant les composantes supersingulières de $K$, on trouve dans l'image de
\[ \Pi_1(\widetilde{\mathcal X})\longrightarrow\mathbb D_{2n} \]
les éléments
\[ t\longrightarrow\lambda'(t)=t-2(n-1). \]
Or le pgcd de $2(n-1)$ et $n$ est tjrs 2. Donc l'image contient tjrs $\mathbb Z/n\mathbb Z$, d'indice 4}

(le quotient étant $\mathbb D_2\simeq\mathbb F_2\times\mathbb F_2$)

\deleted{Ceci ne nous dit \inserted{rien} sur l'image de $\Pi_1(\mathcal X)$ \unsure{ou} $\Pi_1(\widetilde{\mathcal X})$ dans $\mathbb D_2$, et en particulier sur son image dans $\pm1$, \inserted{par la} \unsure{composée|composé} avec
\[ \mathbb D_{2n}\xrightarrow{\chi}\pm1 \]}

\deleted{On a l'impression que pour avoir des informations, il faut faire pivoter un p.r. autour d'un sommet de $\underline{\Sigma}$ et faire ainsi un cycle et voir l'effet produit. On va faire un demi-tour, à partir d'une position supersingulière. disons, \deleted{\unsure{par}} transport parallèle l'orientation de la position supersingulière aura changé, et \deleted{\unsure{le}} $\chi(\lambda)$ sera égal à $+1$ si le nb total de ``pas'' pour revenir au pt de départ est \uline{impair}, à $-1$ dans le cas contraire. Donc on aimerait pouvoir le faire en un nb \uline{pair} de pas !}

% slanted
\margin{4.1. Attention. ces possibilités sont exclues dans le jeu des Pol(D)}

% slanted
\margin{\unsure{Gribul.|Grimul.} cf notes du 4.1.}
